\section{Estimación del Equipamiento}

La estimación del equipamiento parte de una condición favorable del proyecto: la
infraestructura de captura ya existe en el recinto y se reutiliza. El sistema no
requiere la instalación de cámaras nuevas, sino que consume el flujo de video
producido por el grabador digital que opera el parqueadero. En consecuencia, la
inversión se concentra en la capacidad de cómputo necesaria para ejecutar la
inferencia y en el almacenamiento requerido para el material de video y los
conjuntos de datos derivados.

La Tabla \ref{tab:equipamiento} enumera los elementos considerados. Se distingue
entre los elementos que representan un desembolso atribuible al proyecto y
aquellos que se reutilizan de la infraestructura existente y que, por tanto, no
se computan como inversión.

\begin{longtable}{
        P{0.25}
        P{0.35}
        P{0.12}
        P{0.18}
    }
    \caption{Equipamiento requerido por el proyecto}
    \label{tab:equipamiento}\\

    \toprule
    \rowcolor{gray!20}
    \textbf{Elemento} &
    \textbf{Especificación} &
    \textbf{Cantidad} &
    \textbf{Costo} \\
    \midrule
    \endfirsthead

    \multicolumn{4}{c}{\textit{Continuación de la Tabla
    \ref{tab:equipamiento}}} \\
    \toprule
    \rowcolor{gray!20}
    \textbf{Elemento} &
    \textbf{Especificación} &
    \textbf{Cantidad} &
    \textbf{Costo} \\
    \midrule
    \endhead

    Equipo de cómputo con acelerador gráfico &
    GPU NVIDIA GeForce RTX 3050 (4 GB de memoria), procesador de 12 núcleos,
    15 GiB de memoria principal &
    1 & Bs 10.000 \\
    \midrule

    Almacenamiento externo &
    Unidad de $1{,}9$ TB destinada al material de video y a los conjuntos de datos &
    1 & Bs 850 \\
    \midrule

    Infraestructura de captura &
    Grabador digital y cámaras de los puntos de acceso, preexistentes en el
    recinto &
    --- & Sin desembolso \\
    \midrule

    Red y energía &
    Infraestructura de red local y suministro eléctrico del recinto &
    --- & Sin desembolso \\
    \midrule

    Licencias de software &
    Totalidad del componente de software bajo licencias de código abierto &
    --- & Sin desembolso \\
    \midrule

    \multicolumn{3}{r}{\textbf{Total de equipamiento}} & \textbf{Bs 10.850} \\*
    \multicolumn{4}{r}{\footnotesize\textit{Fuente: Elaboración propia}} \\*
    \bottomrule
\end{longtable}

El dimensionamiento del almacenamiento se deriva del volumen efectivo del
material procesado. La captura del recinto se realiza a una resolución de
$1280 \times 720$ píxeles, a una tasa de quince fotogramas por segundo y con
codificación de alta eficiencia, lo que produce aproximadamente 6,7 GB por cada
treinta horas de registro. La generación de los conjuntos de datos añade, a su
vez, poco más de 3 GB derivados del mismo material. Estas cifras sitúan los
requerimientos de almacenamiento muy por debajo de la capacidad disponible, pero
justifican su consideración explícita: el proceso de extracción recorre
repetidamente el archivo de video y requiere conservar simultáneamente las
muestras emitidas.

El equipo de cómputo merece una precisión. El acelerador gráfico empleado dispone
de apenas 4 GB de memoria, un orden de magnitud inferior al de los equipos
habitualmente asociados al entrenamiento de modelos de detección. Esta
restricción no es un inconveniente del proyecto sino su condición de contorno:
constituye la evidencia de que el sistema puede operar sobre hardware de gama
moderada, que es precisamente el escenario de despliegue para el que se diseñó la
estrategia de optimización descrita en el capítulo anterior. El mismo equipo
ejecuta el servicio de gestión, el trabajador de inferencia y la base de datos
como contenedores, sin requerir un servidor dedicado, lo que elimina una partida
de inversión que el diseño evita por construcción.

Respecto de las licencias, la totalidad de los componentes empleados ---lenguaje,
marco de trabajo web, motor de base de datos, cola de tareas, bibliotecas de
aprendizaje profundo, motor de reconocimiento de texto y biblioteca de interfaz---
se distribuye bajo licencias de código abierto, de modo que no generan costo de
adquisición. Se deja constancia, no obstante, de que la biblioteca de detección
empleada se distribuye bajo una licencia con obligaciones de reciprocidad que
condicionan su uso en un producto comercial cerrado; esta consideración no afecta
al presente proyecto, de carácter académico, pero resulta relevante para una
eventual explotación comercial del sistema.

Los precios unitarios se declaran a la fecha de cierre del documento: Bs 10.000
para el equipo de cómputo y Bs 850 para la unidad de almacenamiento, lo que
arroja un total de Bs 10.850. Al igual que la tarifa horaria, son parámetros de
mercado sujetos a variación y deben reemplazarse por cotización formal antes de la
presentación. La partida es pequeña frente al costo de desarrollo, de modo que una
variación de precios en el equipamiento no altera la conclusión del análisis
económico.
